\documentclass[11pt]{article}
\usepackage[a4paper,margin=26mm]{geometry}
\usepackage[T1]{fontenc}
\usepackage{lmodern}
\usepackage{amsmath,amssymb,amsthm}
\usepackage{microtype}
\usepackage{array}
\usepackage{xurl}
\usepackage{hyperref}
\hypersetup{colorlinks=true,linkcolor=blue,urlcolor=blue,citecolor=blue,
 pdftitle={A terrace for the alternating group A7},
 pdfauthor={Horace Dong}}
\newtheorem{theorem}{Theorem}
\newtheorem{lemma}{Lemma}
\newtheorem{proposition}{Proposition}
\newtheorem{corollary}{Corollary}
\theoremstyle{definition}
\newtheorem{definition}{Definition}
\theoremstyle{remark}
\newtheorem{remark}{Remark}
\newcommand{\repoversion}{1.25.0}
\setlength{\emergencystretch}{1.5em}
\newcommand{\Alt}{A_7}
\newcommand{\inv}{^{-1}}

\title{A terrace for the alternating group $A_7$}
\author{Horace Dong\\ \small Independent researcher\\
\small\texttt{maxwelldhx@gmail.com}\\
\small ORCID: \href{https://orcid.org/0009-0002-7554-7548}{0009-0002-7554-7548}}
\date{22 November 2026}
\begin{document}
\maketitle

\begin{abstract}
Let $G$ be a finite group of order $n$. An arrangement $a_1,\dots,a_n$ of the
elements of $G$ is a terrace if, among the $n-1$ quotients
$a_i^{-1}a_{i+1}$, every involution occurs exactly once and every other
non-identity element $g$ occurs, together with $g^{-1}$, exactly twice in
total. Bailey conjectured that every finite group except the non-cyclic
elementary abelian $2$-groups has a terrace. We give an explicit terrace of
the alternating group $A_7$, of order $2520$. Apart from an asymptotic theorem
without an explicit bound, none of the families of terraced groups listed in Ollis's 2025 survey is known to contain $A_7$, so this is a further verified case of Bailey's
conjecture. Combined with a theorem of Anderson and Ihrig, the terrace also
gives symmetric sequencings for every group $G$ with a unique involution such
that $G/O(G)\Lambda(G)\cong A_7$, where $O(G)$ is the largest normal subgroup
of odd order and $\Lambda(G)$ is generated by the involution; this is the
$A_7$ case of their reduction of Anderson's conjecture on symmetric
sequencings. The terrace
is not directed, so it says nothing about whether $A_7$ is sequenceable. It
was found by a hill-climbing search, and the fact that it is a terrace is
formally verified in Lean~4 by a checking program that is proved correct and
evaluated by the Lean kernel. The search and the formalization were carried
out by AI agents working under the author's direction; their role is described
at the end of the paper.
\end{abstract}

\section{Introduction}

Let $G$ be a finite group of order $n$ with identity $e$. Following
Ollis~\cite[Section~1]{ollis16}, for an arrangement $a=(a_1,\dots,a_n)$ of the
elements of $G$ put
\[
b_i=a_i^{-1}a_{i+1}\qquad(1\le i\le n-1).
\]
If every involution of $G$ occurs exactly once in $b=(b_1,\dots,b_{n-1})$ and,
for every $x\in G$ with $x^2\ne e$, the elements $x$ and $x^{-1}$ occur
``twice in total'' in $b$, then $a$ is a \emph{terrace} for $G$ and $b$ is its
\emph{$2$-sequencing}; a group with a terrace is \emph{terraced}. The two
occurrences of a pair $\{x,x^{-1}\}$ may be $x$ twice, $x^{-1}$ twice, or one
of each. The identity never occurs in $b$, since the $a_i$ are distinct.
Left-multiplying every $a_i$ by $a_1^{-1}$ does not change $b$, so a terraced
group has a \emph{basic} terrace, one with $a_1=e$~\cite[p.~2]{ollis16}. If
the $b_i$ are all distinct, $a$ is a \emph{directed terrace} and $b$ a
\emph{sequencing}, and a group with a sequencing is
\emph{sequenceable}~\cite[Section~1]{ollis16}. A directed terrace is a terrace.
The dynamic survey of Ollis~\cite[Sections~2.3 and~3.4]{survey} gives the
same definition in the form of $2$-sequencings $(e,d_2,\dots,d_n)$ and notes
that terraces and $2$-sequencings are equivalent.

Ollis attributes the notion of a terrace and the terminology to
Bailey~\cite{bailey}, who used terraces to construct quasi-complete Latin
squares; we have not read Bailey's paper and rely on the two texts of
Ollis for the definition. The non-cyclic elementary abelian $2$-groups have
no terrace: all their non-identity elements are involutions, so a terrace
would be a sequencing, and the product $b_1\cdots b_{n-1}$ of all
non-identity elements is $e$, which would force $a_n=a_1$. Bailey's
conjecture~\cite[Conjecture~33]{survey} is that these are the only
exceptions:

\begin{quote}
\textbf{Bailey's conjecture.} All finite groups, except the elementary
abelian $2$-groups of order at least~$4$, are terraced.
\end{quote}

\noindent It is proved for abelian groups~\cite[Theorem~31]{survey}; the
families of non-abelian groups known to be terraced are collected in
Theorem~32 of the survey and recalled in Section~\ref{sec:context}. Our
result concerns the alternating group $\Alt$, which we realize as the group
of even permutations of $\{0,1,\dots,6\}$ with the product
$(pq)(x)=p(q(x))$. It has $7!/2=2520$ elements.

\begin{theorem}\label{thm:main}
There is a bijection $i\mapsto a_i$ from $\{0,1,\dots,2519\}$ onto $\Alt$
with $a_0=e$ such that, writing $d_i=a_i^{-1}a_{i+1}$ for $0\le i\le 2518$,
for every $g\in\Alt$ with $g\ne e$
\[
\#\{\,i : 0\le i\le 2518,\ d_i=g\ \text{or}\ d_i=g^{-1}\,\}=
\begin{cases}1&\text{if } g=g^{-1},\\ 2&\text{otherwise.}\end{cases}
\]
In other words, $(a_0,\dots,a_{2519})$ is a basic terrace for $\Alt$, and
$\Alt$ is terraced.
\end{theorem}

\noindent We index from $0$, as the Lean files do; $d_i$ is the quotient
$b_{i+1}$ in the notation above. For an involution $g$ the condition says that
$g$ occurs exactly once, and for any other $g\ne e$ that $g$ and $g^{-1}$
occur twice in total, so the theorem states exactly the definition of a basic
terrace. The group $\Alt$ has $105$ involutions (the products of two disjoint
transpositions, $3\binom74=105$ of them) and the remaining $2414$
non-identity elements form $1207$ pairs $\{g,g^{-1}\}$; the required
multiplicities add up to $105+2\cdot1207=2519$, the number of quotients.

Theorem~\ref{thm:main} is formally verified in Lean~4
(Section~\ref{sec:lean}). The terrace is explicit: its $2520$ elements are
stored in the Lean file \texttt{Data.lean} of the repository~\cite{repo}, in
a form described in Section~\ref{sec:search}, and we do not print them here.

\paragraph{Why $A_7$.}
None of the families of Theorem~32 of the survey is known to contain $\Alt$
(Section~\ref{sec:context}): $\Alt$ is simple, its order $2520$ exceeds the
range of the computer searches, which reach order $511$, and it is not known
to be sequenceable. Müyesser and Pokrovskiy~\cite[Theorem~6.12]{mp} proved
that every sufficiently large non-abelian group is sequenceable, hence
terraced, but they give no explicit bound; the survey reports an estimate of
Müyesser that groups of order at least about $10^{50}$ would comfortably
suffice for the argument as written~\cite[p.~18]{survey}. So, as far as our
searches show, Theorem~\ref{thm:main} adds one verified group to Bailey's
conjecture. The conjecture itself remains open.

The group $\Alt$ also appears by name in a second problem. A \emph{binary}
group is a group with exactly one involution, and Anderson's conjecture is
that the quaternion group $Q_8$ is the only binary group without a symmetric
sequencing (a sequencing of a particular symmetric form)~\cite[Conjecture~16
and Section~2.3]{survey}. According to Ollis's account~\cite[Section~3]{ollis16}
and the survey~\cite[p.~15]{survey}, Anderson and Ihrig~\cite{ai93b} reduced
the insoluble case to finding $2$-sequencings (equivalently, terraces) of
$\Alt$, $\mathrm{PSL}(2,q)$ and $\mathrm{PGL}(2,q)$ for odd prime powers
$q>3$. The survey lists such terraces for $\mathrm{PSL}(2,5)\cong A_5$,
$\mathrm{PSL}(2,q)$ with $q\in\{7,9,11\}$ and $\mathrm{PGL}(2,q)$ with
$q\in\{5,7\}$~\cite[p.~15]{survey}, and none for $\Alt$. Through Theorem~9 of
the survey, our terrace gives symmetric sequencings for every binary group $G$
with $G/O(G)\Lambda(G)\cong\Alt$ (Corollary~\ref{cor:anderson}). This settles
neither Anderson's conjecture nor the general cases $\mathrm{PSL}(2,q)$ and
$\mathrm{PGL}(2,q)$.

Finally, the terrace of Theorem~\ref{thm:main} is not directed: some
quotients occur twice (Remark~\ref{rem:notdirected}). Whether $\Alt$ is
sequenceable, as Keedwell's conjecture~\cite[Conjecture~15]{survey} predicts,
remains open as far as we know.

\subsection*{Prior work}
The hill-climbing method used to find the terrace, which reverses segments of
an arrangement, is that of Ollis~\cite[Section~2]{ollis16}, who builds on
Anderson's algorithm for directed terraces. With this method and its
directed version without reversals, together with
\cite[Theorem~2.4]{ollis16}, Ollis showed that all non-abelian groups of order at most $511$, except possibly at orders $256$
and $384$, are terraced, and that all non-abelian groups of order $10$ to
$255$ and the groups $A_6$, $S_6$ and some small projective groups are
sequenceable~\cite[Theorems~2.1 and~2.5]{ollis16}; Anderson had earlier
sequenced $A_5$ and $S_5$. A dataset of Yang~\cite{yang}, deposited in
September 2026, reports terraces for all non-abelian groups of orders $256$
and $384$. None of these covers $\Alt$, and searches up to 20 November 2026
found no terrace, directed terrace or sequencing of $\Alt$ and no
non-asymptotic theorem that covers it; the searches and their gaps are
described in the final section.

\section{How the terrace was found and how to check it}\label{sec:search}

\paragraph{The search.}
The terrace was found by a short C++ program written by Codex, a hill-climbing
search in the style of Ollis~\cite[Section~2]{ollis16}. The program lists the
$2520$ even permutations, tabulates for every ordered pair $(u,v)$ the pair
$\{w,w^{-1}\}$ with $w=u^{-1}v$, and starts from a random arrangement (a
shuffle with the Mersenne Twister generator and seed $1993$). For each
non-identity pair $\{g,g^{-1}\}$ its \emph{quota} is $1$ if $g$ is an
involution and $2$ otherwise, and the \emph{excess} of an arrangement is the
sum, over all pairs, of the amount by which the number of quotients in the
pair exceeds the quota. An arrangement is a terrace exactly when its excess
is~$0$.

The only move is to reverse an initial segment $(a_1,\dots,a_{k+1})$. Inside
the segment each quotient is replaced by its inverse (and the order of the
quotients is reversed), so the pairs of these quotients do not change; the
only quotient that changes is the one at the seam, $a_{k+1}^{-1}a_{k+2}$,
which becomes $a_1^{-1}a_{k+2}$. In each step the program reverses the whole
arrangement with probability $1/2$ (this keeps the excess) and then picks
at random a seam move that lowers the excess by one if there is such a move,
and otherwise a move that leaves the excess unchanged. Its log reports an
excess of $680$ after one step and excess $0$ after $1544$ steps. The random generator is used
only to find the terrace; the shuffle of a different C++ standard library may
lead to a different run, but the terrace found is stored and its checks use no
randomness.

\paragraph{The stored terrace.}
Codex's program printed the terrace as $2520$ lines of seven digits, the images
of $0,\dots,6$. A deterministic Python encoder left-multiplied every element
by the inverse of the initial one, which gives a basic terrace with the same
quotients, and wrote the result, together with auxiliary tables, to
\texttt{Data.lean}; rerunning the encoder reproduced \texttt{Data.lean} and
\texttt{Cert.lean} byte for byte. In \texttt{Data.lean} every table is a function of $i$ given by
hexadecimal constants $B_0,B_1,\dots$ ($40$ of them for the tables indexed by
the $2520$ rows), each packing $64$ consecutive entries of a fixed width~$w$:
\[
\mathrm{table}(i)=\bigl\lfloor B_{\lfloor i/64\rfloor}\big/2^{\,w\,(i\bmod 64)}\bigr\rfloor\bmod 2^{w}.
\]
For the table \texttt{images} ($w=21$), the element $a_i$ is the permutation
\[
a_i(k)=\bigl\lfloor \mathtt{images}(i)/8^{k}\bigr\rfloor\bmod 8\qquad(0\le k\le 6).
\]
For example $a_0=e$, and in image notation $[a(0),\dots,a(6)]$ the next
elements are $a_1=[5,6,1,3,0,4,2]$ and $a_2=[0,5,2,4,3,1,6]$; the last is
$a_{2519}=[4,6,0,2,1,3,5]$.

\paragraph{Checking it without Lean.}
Theorem~\ref{thm:main} can be checked directly from this table by a short
program: decode the $2520$ rows; check that they are pairwise distinct even
permutations (hence all of $\Alt$) and that row $0$ is the identity; compute
the $2519$ quotients $d_i=a_i^{-1}a_{i+1}$ with the product above; and count,
for every $g\ne e$, the indices $i$ with $d_i\in\{g,g^{-1}\}$. The terrace was
checked in this way by two Python programs written by Codex on the original
witness, by the final reviewer on \texttt{Data.lean}, and by Claude Code
while writing this paper (see the final section). Such checks are ordinary
computations; the Lean proof described next does not depend on them.

\begin{remark}\label{rem:notdirected}
The terrace is not directed. With the numbering of Theorem~\ref{thm:main},
$d_3=d_{36}=(0\,3\,1)(4\,6\,5)$, whose image list is $[3,0,2,1,6,4,5]$. By
our count, $592$ of the $1207$ pairs $\{g,g^{-1}\}$ occur as the same element
twice and $615$ as $g$ and $g^{-1}$ once each, so only $1927$ different
quotients occur. These facts are computations outside Lean.
\end{remark}

\section{The certificate checked in Lean}\label{sec:cert}

A direct formalization of the check above would ask the Lean kernel to
compare $2519$ quotients with one another. The Lean proof instead checks a
certificate whose conditions are local to one row or one quotient, and derives
the counting statement by a short abstract argument. We describe both parts.

\begin{lemma}[counting from a matching]\label{lem:matching}
Let $G$ be a finite group of order $n\ge2$, let $i\mapsto a_i$ be a bijection
from $\{0,\dots,n-1\}$ onto $G$ with $a_0=e$, and put
$d_i=a_i^{-1}a_{i+1}$ for $0\le i\le n-2$. Suppose there is a permutation
$\sigma$ of $\{0,\dots,n-2\}$ such that, with $r_i=a_{\sigma(i)+1}$,
\[
d_i=r_i\quad\text{or}\quad d_i=r_i^{-1}\qquad\text{for every } i .
\]
Then $(a_0,\dots,a_{n-1})$ satisfies the conclusion of
Theorem~\ref{thm:main} (with $n-2$ in place of $2518$): for every $g\ne e$,
the number of $i$ with $d_i\in\{g,g^{-1}\}$ is $|\{g,g^{-1}\}|$.
\end{lemma}

\begin{proof}
Fix $g\ne e$ and let $S=\{g,g^{-1}\}$. Since $d_i\in\{r_i,r_i^{-1}\}$ and $S$
is closed under inversion, $d_i\in S$ if and only if $r_i\in S$. So
$i\mapsto r_i$ maps $I=\{i:d_i\in S\}$ into $S$. It is injective, because
$\sigma$, $j\mapsto j+1$ and $j\mapsto a_j$ are injective. It is surjective:
an element $x\in S$ is not $e$, so $x=a_m$ with $m\ne0$, that is $m=j+1$ for
some $j$, and $i=\sigma^{-1}(j)$ satisfies $r_i=x$, hence $i\in I$. Therefore
$|I|=|S|$, which is $1$ if $g=g^{-1}$ and $2$ otherwise.
\end{proof}

\paragraph{The certificate.}
The tables in \texttt{Data.lean} provide, for each row $i$, the packed image
list $\mathtt{images}(i)$, a word $\mathtt{words}(i)$ and, for each $i<2519$,
an index $\mathtt{edge}(i)<2519$ with a direction bit, together with a reverse
table \texttt{edgeBack} and the tables of a small hash function. Nothing about
these tables is trusted; they are only data for the following checks.
\begin{enumerate}
\item[(R1)] \emph{Rows are genuine elements.} The word $\mathtt{words}(i)$ is
read as four base-$36$ digits; each digit gives a product
$(0\ j)(0\ j')$ of two transpositions with $j,j'\in\{1,\dots,6\}$, which is
an even permutation. The product $v_i$ of the four factors is therefore an
element of $\Alt$ by construction, and the check is that $v_i(k)$ equals the
stored image $a_i(k)$ for $k=0,\dots,6$. (The encoder found for each row a
shortest product of star transpositions $(0\ j)$ by breadth-first search;
these products have even length at most eight, and shorter words are padded
with the trivial pair $(0\ 1)(0\ 1)$.)
\item[(R2)] \emph{Rows are distinct.} The seven images of row $i$ are packed
into the integer $\sum_k a_i(k)\,8^k$, and an explicit function
$\mathrm{rank}$ (a lookup in a table of size $4093$ after a shift that depends
on the residue modulo $257$) must return $i$ on it. Hence two rows with the
same images have the same index.
\item[(S)] \emph{Quotients are matched.} For $i<2519$ let $s=\mathtt{edge}(i)$.
The check is that $\mathtt{edgeBack}(s)=i$ and, according to the direction
bit, either $a_i\circ a_{s+1}=a_{i+1}$ or $a_{i+1}\circ a_{s+1}=a_i$, compared
point by point on $\{0,\dots,6\}$.
\end{enumerate}
These checks are combined into one Boolean function of $i$, and the kernel
evaluates it on $79$ blocks of at most $32$ consecutive rows
($78$ blocks of $32$ and one of $24$), one \texttt{decide +kernel} call per
block. A lemma proved by induction turns each evaluation into the statement
that the check holds on the block, and a second lemma merges adjacent blocks,
giving the check for all $i<2520$. A wrong or colliding table entry would make
the evaluation of its block fail.

\begin{proof}[Proof of Theorem~\ref{thm:main}]
By (R1), $i\mapsto v_i$ maps $\{0,\dots,2519\}$ into $\Alt$, and $v_i$ acts
on $\{0,\dots,6\}$ as the stored row $a_i$; we identify $a_i$ with $v_i$. By
(R2) this map is injective, and since $|\Alt|=2520$ it is a bijection. Row $0$
is the identity, so $a_0=e$. By (S), $i\mapsto\mathtt{edge}(i)$ is injective
on $\{0,\dots,2518\}$ (it has the left inverse \texttt{edgeBack}), hence a
permutation $\sigma$. If the direction bit of $i$ is positive, then
$a_{i+1}=a_i a_{\sigma(i)+1}$, so $d_i=a_{\sigma(i)+1}$; otherwise
$a_i=a_{i+1}a_{\sigma(i)+1}$, so $d_i=a_{\sigma(i)+1}^{-1}$. Lemma~\ref{lem:matching}
now gives the theorem.
\end{proof}

\noindent The matching $\sigma$ pairs each quotient $d_i$ with a vertex
$a_{\sigma(i)+1}\ne e$ that equals $d_i$ or $d_i^{-1}$. It exists because
$i\mapsto a_{i+1}$ runs through $\Alt\setminus\{e\}$, and each pair
$\{g,g^{-1}\}$ contains exactly as many vertices as the terrace has quotients
in it. The encoder found $\sigma$ greedily from the terrace.

\section{Consequences}\label{sec:context}

The statements of this section are not formalized; they combine
Theorem~\ref{thm:main} with published theorems.

\paragraph{Bailey's conjecture.}
Theorem~32 of the survey~\cite{survey} states that the following non-abelian
groups are terraced: (i)~sequenceable groups; (ii)~groups of odd order;
(iii)~groups with a terraced normal subgroup of odd index, and groups with a
normal subgroup of odd order and a terraced quotient (see
also~\cite[Theorem~2.4]{ollis16}); (iv)~the semidihedral groups $SD_{8t}$
and the groups $M_{8t}$, $t\ge2$; (v)~various groups with a central subgroup
isomorphic to $C_2^2$; (vi)~all groups of order at most $511$, except possibly
at orders $256$ and $384$. The group $\Alt$ is in none of these families as
far as is known. It is not known to be sequenceable: the survey's list of
groups known to be sequenceable~\cite[Section~2.6]{survey} contains $A_5$,
$S_5$, $A_6$, $S_6$ and some projective groups, and all non-abelian groups of order $10$ to $255$, but not $\Alt$. It has even order, so (ii) does not
apply. Being simple, it has
no proper non-trivial normal subgroup from which (iii) could lift a terrace.
Its centre is trivial and it has no cyclic subgroup of index~$2$, so (iv) and
(v) do not apply, and $2520>511$. Theorem~14 of the survey (Müyesser and
Pokrovskiy~\cite[Theorem~6.12]{mp}) covers all sufficiently large non-abelian
groups, with no explicit bound. In this sense Theorem~\ref{thm:main} is a
further verified case of Bailey's conjecture, as far as the searches up to
20 November 2026 show.

Combined with (iii), Theorem~\ref{thm:main} also shows that every group with
a normal subgroup $N$ of odd order and $G/N\cong\Alt$, and every group
containing $\Alt$ as a normal subgroup of odd index, is terraced; for
example $\Alt\times H$ for every group $H$ of odd order. Groups of this kind
that are large enough are already covered by the theorem of Müyesser and
Pokrovskiy.

\paragraph{Anderson's conjecture.}
For a binary group $G$ let $\Lambda(G)$ be the subgroup generated by its
involution, which is central, and let $O(G)$ be the largest normal subgroup of
odd order. The survey states the following two results.

\begin{itemize}
\item \cite[Theorem~8]{survey} A binary group $G$ of order $2n$ has a symmetric
sequencing if and only if $G/\Lambda(G)$ has a $2$-sequencing.
\item \cite[Theorem~9]{survey}, attributed to Anderson and Ihrig~\cite{ai93a}:
if $G$ is a binary group and $G/O(G)\Lambda(G)$ has a $2$-sequencing, then $G$
has a symmetric sequencing.
\end{itemize}

\begin{corollary}\label{cor:anderson}
Every binary group $G$ with $G/O(G)\Lambda(G)\cong\Alt$ has a symmetric
sequencing.
\end{corollary}

\begin{proof}
By Theorem~\ref{thm:main}, $(e,d_0,d_1,\dots,d_{2518})$ is a $2$-sequencing of
$\Alt$ in the sense of~\cite[p.~14]{survey}: its partial products are the basic
terrace $a_0,\dots,a_{2519}$. Apply Theorem~9 of the survey.
\end{proof}

\noindent This is the $A_7$ part of the reduction mentioned in the
introduction. The reduction itself, in which the group $\Alt$ appears as the
middle factor $G_1/G_2$ of a subnormal series $\{e\}\trianglelefteq
G_2\trianglelefteq G_1\trianglelefteq G/\Lambda(G)$ whose other factors have odd
order, and the statement that each of the groups $\Alt$, $\mathrm{PSL}(2,q)$,
$\mathrm{PGL}(2,q)$ accounts for infinitely many insoluble binary groups, are
due to Anderson and Ihrig~\cite{ai93b}; we have read them only in Ollis's
paraphrase~\cite[Section~3]{ollis16} and in the survey~\cite[p.~15]{survey},
not in the original paper.

Without the terrace, part of Corollary~\ref{cor:anderson} already followed
from published results. If $G/O(G)\Lambda(G)\cong\Alt$, then $G/\Lambda(G)$
maps onto $\Alt$ and is therefore non-abelian, of order $|G|/2$. When $|G|$ is
large enough, $G/\Lambda(G)$ is sequenceable by the theorem of Müyesser and
Pokrovskiy, a sequencing is a $2$-sequencing~\cite[p.~14]{survey}, and
Theorem~8 gives a symmetric sequencing of $G$. Since no explicit bound is
known (the survey quotes an estimate of about $10^{50}$), this covers the
binary groups of this branch only above an unknown order. With the terrace,
Theorem~9 covers every $G$ with $G/O(G)\Lambda(G)\cong\Alt$, whatever its
order.

\paragraph{Keedwell's conjecture.}
Keedwell conjectured that $D_6$, $D_8$ and $Q_8$ are the only non-abelian
groups that are not sequenceable~\cite[Conjecture~15]{survey}. For $\Alt$ this
needs a directed terrace, and the terrace of Theorem~\ref{thm:main} is not
directed (Remark~\ref{rem:notdirected}). Large non-abelian groups are
sequenceable by~\cite[Theorem~6.12]{mp}, but for $\Alt$ the question is open
as far as we know.

{\sloppy
\section{Formal verification}\label{sec:lean}

Theorem~\ref{thm:main} is formally verified in Lean~4 (version 4.34.1) with
Mathlib~\cite{mathlib} (commit \texttt{d13f23b7}). The $8$ files are in
version~\repoversion{} of the repository~\cite{repo}, directory
\texttt{lean/Research/TaskR/A7/}. Lemma~\ref{lem:matching} is proved there for
the group $\Alt$ and $n=2520$. Remark~\ref{rem:notdirected} and
Section~\ref{sec:context} are not part of the Lean files.

\paragraph{The statement.}
The statement file \texttt{Challenge.lean} contains only definitions and the
statement, no proofs. In ASCII transliteration (docstrings and the
\texttt{import}, \texttt{namespace} and \texttt{end} lines omitted; the file
itself uses Unicode, e.g.\ $\simeq$, $\exists$, $\forall$, $\ne$ and
${}^{-1}$):
{\footnotesize
\begin{verbatim}
abbrev G := alternatingGroup (Fin 7)
def difference (a : Fin 2520 -> G) (i : Fin 2519) : G :=
  (a i.castSucc)^-1 * a i.succ
def HasTerrace : Prop :=
  exists a : Fin 2520 ~= G, a 0 = 1 /\
    forall g : G, g != 1 ->
      Fintype.card {i : Fin 2519 // difference a i = g \/ difference a i = g^-1} =
        if g = g^-1 then 1 else 2
\end{verbatim}}
\noindent Here \texttt{alternatingGroup (Fin 7)} is Mathlib's alternating
group on $\{0,\dots,6\}$, the subgroup of even permutations of the
permutation group, in which $(pq)(x)=p(q(x))$; \texttt{Fin 2520 \textasciitilde= G}
(written $\simeq$ in the file) is the type of bijections from
$\{0,\dots,2519\}$ onto $G$; \texttt{1} is the identity; for
$i\in\{0,\dots,2518\}$, \texttt{i.castSucc} and \texttt{i.succ} are $i$ and
$i+1$ as elements of $\{0,\dots,2519\}$, so \texttt{difference a i} is
$d_i=a_i^{-1}a_{i+1}$; \texttt{!=} stands for $\ne$; and
\texttt{Fintype.card \{i : Fin 2519 // \dots\}} is the number of indices $i$
with $d_i=g$ or $d_i=g^{-1}$. So \texttt{HasTerrace} is
Theorem~\ref{thm:main}.

The final reviewer (see the last section) checked the following points.
(i)~A bijection lists every element of $\Alt$ exactly once. (ii)~The identity
is excluded from the count, but it cannot occur as a quotient anyway.
(iii)~The condition \texttt{a 0 = 1} asks for a basic terrace; this makes the
statement stronger, and is equivalent to the existence of a terrace by left
multiplication. (iv)~The quotients $a_i^{-1}a_{i+1}$ follow Ollis's
convention. Under the other convention $a_{i+1}a_i^{-1}$ this particular
arrangement is not a terrace, but the two conventions are exchanged by
replacing every $a_i$ by $a_i^{-1}$, so existence is the same; for the same
reason the order of composition in Mathlib does not matter. (v)~In an audit
file of its own, the reviewer checked that \texttt{G} is
\texttt{alternatingGroup (Fin 7)} by definition, that this group has $2520$
elements, and that the fully unfolded statement, written with Mathlib names
only, follows from the final theorem. (vi)~The statement has no extra
hypotheses, and the source shows that the property is not trivial (the Klein
four-group has no terrace).

\paragraph{Review and freezing.}
\texttt{Challenge.lean} was reviewed against~\cite{ollis16} and~\cite{survey}
by a Codex agent that had not written it, and frozen, by recording its
SHA-256 hash, on 20 November 2026 at 19:32~UTC. The development logs record
only a compilation of the statement file itself before that time, and the
file was not changed afterwards. The terrace had been found before (see the
last section); the statement does not contain it.

\paragraph{The proofs.}
In \texttt{Proof.lean}, \texttt{Proof.HasTerrace} is \texttt{fromMatching}
(Lemma~\ref{lem:matching}, file \texttt{Counting.lean}) applied to
\texttt{Interpret.ordering}, \texttt{Interpret.first},
\texttt{Interpret.matching} and \texttt{Interpret.edges}, each instantiated
with \texttt{Cert.allRows}. Table~\ref{tab:lean} lists the Lean counterparts
of the steps of Section~\ref{sec:cert} (names in the namespace
\texttt{TaskR.A7}).

\begin{table}[ht]
\centering\footnotesize
\renewcommand{\arraystretch}{1.05}
\begin{tabular}{@{}>{\raggedright\arraybackslash}p{0.42\textwidth}>{\raggedright\arraybackslash}p{0.54\textwidth}@{}}
\hline
Step of the written proof & Lean (file) \\
\hline
The tables &
\texttt{Data.images}, \texttt{Data.words}, \texttt{Data.rankTable},
\texttt{Data.shifts}, \texttt{Data.hashMod}, \texttt{Data.edgeTags},
\texttt{Data.edgeBack} (\texttt{Data}) \\
(R1): star-transposition words give elements of $\Alt$ &
\texttt{Kernel.point}, \texttt{Kernel.pair} (via Mathlib's
\texttt{Equiv.\allowbreak Perm.\allowbreak mul\_\allowbreak mem\_\allowbreak
alternatingGroup\_\allowbreak of\_\allowbreak isSwap}),
\texttt{Kernel.word}, \texttt{Kernel.vertex} (\texttt{Kernel}) \\
Decoding and packing a row; the rank function &
\texttt{Kernel.raw}, \texttt{Kernel.pack}, \texttt{Kernel.rank}
(\texttt{Kernel}) \\
(S): the matching and its direction bit &
\texttt{Kernel.edge}, \texttt{Kernel.positive} (\texttt{Kernel}) \\
The Boolean checks (R1), (R2), (S) &
\texttt{Kernel.rowOK}, \texttt{Kernel.stepOK}, \texttt{Kernel.valid}
(\texttt{Kernel}) \\
Block evaluation and merging &
\texttt{Kernel.interval}, \texttt{Kernel.Good},
\texttt{Kernel.interval\_sound}, \texttt{Kernel.merge} (\texttt{Kernel});
$79$ private block theorems from \texttt{b0\_32} to \texttt{b2496\_24} and
\texttt{Cert.allRows} (\texttt{Cert}) \\
The rows form a bijection onto $\Alt$ with $a_0=e$ &
\texttt{Interpret.ordering} (using Mathlib's
\texttt{card\_\allowbreak alternatingGroup}), \texttt{Interpret.first}
(\texttt{Interpret}) \\
$\sigma$ is a permutation; $d_i\in\{r_i,r_i^{-1}\}$ &
\texttt{Interpret.matching}, \texttt{Interpret.edges} (\texttt{Interpret}) \\
Lemma~\ref{lem:matching} for $\Alt$ &
\texttt{fromMatching} (\texttt{Counting}) \\
Theorem~\ref{thm:main} &
\texttt{Proof.HasTerrace} (\texttt{Proof}),
\texttt{check\_HasTerrace} (\texttt{CheckA7}) \\
\hline
\end{tabular}
\caption{Lean counterparts of the steps of the proof.}
\label{tab:lean}
\end{table}

\noindent The test file \texttt{CheckA7.lean} restates the theorem with the
frozen type (\texttt{check\_\allowbreak HasTerrace}). The axioms of
\texttt{check\_\allowbreak HasTerrace} and of \texttt{Proof.HasTerrace},
printed by \texttt{CheckA7.lean}, are exactly \texttt{propext},
\texttt{Classical.choice} and \texttt{Quot.sound}; the final reviewer found
the same for \texttt{Cert.allRows} and \texttt{fromMatching}. No file uses
\texttt{sorry}, \texttt{native\_\allowbreak decide},
\texttt{bv\_\allowbreak decide} or an \texttt{axiom} declaration. Apart from
the $79$ calls of \texttt{decide +kernel} in \texttt{Cert.lean}, which have the
kernel evaluate the checks, the \texttt{decide} tactic is used only for small arithmetic facts (such as
$0<6$ and $7!/2=2520$); the Boolean \texttt{decide} inside
\texttt{Kernel.rowOK} and \texttt{Kernel.stepOK} is part of the checks. \texttt{Cert.lean} raises the recursion
depth limit and switches off the heartbeat limit, and
\texttt{Interpret.lean} raises the heartbeat limit; these options affect
resources only.

\paragraph{Checks.}
The $8$ files ($5$ proof modules, \texttt{Challenge.lean},
\texttt{Proof.lean} and \texttt{CheckA7.lean}) were compiled module by module
in two new empty directories on the author's Mac, with the pinned toolchain
and with warnings treated as errors, on 20 November 2026
(19:36--19:40~UTC), and every module was replayed with \texttt{leanchecker};
each directory also ran one build of the whole package. All $34$ commands
($17$ per directory) ended with exit code~$0$, the \texttt{leanchecker} runs
printed nothing, and the two directories produced identical \texttt{.olean}
files. The largest module, \texttt{Cert.lean}, took about $18$~seconds to
compile, with a peak memory use of about $3.2$--$3.3$~GB. The final reviewer (see the
next section) compiled the $8$ files again with plain \texttt{lean}, together
with a short audit file of its own, in a third new directory on 20 November
2026, and replayed every module with \texttt{leanchecker}; according to its
file times, the $8$ project modules were compiled and replayed between 19:43
and 19:45~UTC. All runs ended with exit code~$0$ and printed no warnings (an
earlier draft of its own audit file contained an error and was corrected, and
one \texttt{leanchecker} call on that file, which failed to start for lack of
\texttt{lean} on the path, without a Lean verdict, was repeated). Mathlib and
its dependencies were taken from a shared prebuilt cache at the pinned commits
and were not rechecked; the \texttt{leanchecker} replays cover the $8$ project
modules. In the repository the files have module names under
\texttt{Research.TaskR.A7} and differ from the checked files only in their
\texttt{import} lines (\texttt{Challenge.lean} and \texttt{Data.lean} are
unchanged). The repository versions were made by a script that changed only the local \texttt{import} lines, and on 20 November 2026 (19:55--20:07~UTC) they were compiled and replayed with \texttt{leanchecker} again, module by module, in two new empty directories on the author's Mac, one after the other, together with the files of three other papers of version~\repoversion{} ($39$ modules in all); there were no warnings, and the two directories produced identical \texttt{.olean} files and the same axiom lists. In the whole-project rebuild of version~\repoversion{} of the repository (on the author's Windows computer, in two passes on 20 and 21 November 2026), these files were compiled and replayed with \texttt{leanchecker} again in pass~2 (21 November 2026), which built the files not included in pass~1 on top of its results; no module needed a second attempt (details in the repository). The formalization has not been
checked with comparator or with an independent kernel.

\par}

\section*{How this paper was produced}
This work was carried out from 20 to 22 November 2026 (UTC) in a workflow directed by
the author, using OpenAI Codex with the model GPT-6.2-Astra (reasoning level
ultra) and Anthropic Claude Code with the model Claude Opus~5.5. The author asked for a search for hard and significant open problems.
Claude Code wrote this request for Codex at about 18:04~UTC, and from about
18:37~UTC it was carried out in a Codex session separate from the main one,
in which Codex screened $78$ candidate problems by 18:54~UTC. In this
search it wrote the C++ program of Section~\ref{sec:search}, which found the
terrace (the witness file is dated 18:38~UTC), checked the terrace with two separately written Python programs, searched for
prior work from 18:39~UTC and compared the terrace with it between 18:45 and
18:51~UTC. The author chose this problem and one other. In a second
separate Codex session, from 19:29~UTC, Codex agents (same model) checked the
reading of the sources, searched again for prior work, wrote the statement
file, the encoder, the Lean proof and the written proof, and ran the checks.
The reading of the sources was checked, and the statement file reviewed before
it was frozen, by one newly started Codex agent that had taken no part in
finding the terrace and had written neither the statement nor the proofs.
According to the file timestamps and the development logs, the proof files
were compiled between 19:32 and 19:36~UTC, after the freeze. After the checks,
an independent Claude Code agent (model Claude Opus~5.5) that had written none
of the files gave the final review, starting at 19:42~UTC. It read the sources
itself, compared the statement file with them, read the written proof, the
proof files and the raw logs of the checks, checked the terrace in Python from
the tables in \texttt{Data.lean} and against Codex's witness file, and rebuilt
the formalization as described above. Its decision was APPROVE\_FORMAL, with
six minor remarks: five on the written account (one optional), all taken into
account here, and one asking for a further priority check before release.
Claude Code (Claude Opus~5.5) checked the sources and wrote this text from
Codex's manuscript, its search records and the Lean files; it also decoded
\texttt{Data.lean} with a short Python program of its own, which confirmed
Theorem~\ref{thm:main} and the counts in Remark~\ref{rem:notdirected}.
An independent Claude Code agent (model Claude Opus~5.5) that had written none of the files reviewed this paper, and another reviewed the release; their corrections were applied. Throughout, ``we'' refers to this workflow; the reading and
searching described in the paper were done by the agents. Codex agents
searched for prior work on 20 November 2026, in the search for problems and
again before the statement was frozen, and the final reviewer searched again;
they found no terrace, directed terrace or sequencing of $\Alt$ and no
non-asymptotic theorem that covers it (details in the repository). Some
sources could not be reached: the website with Ollis's search programs cited
in~\cite[p.~3]{ollis16}, the full text of~\cite{ai93b}, and some search
services, which refused or limited automated queries. Claude Code searched again on 21 November 2026 (17:54--18:21~UTC; the arXiv, Semantic Scholar, Crossref, DataCite, Zenodo, MathDB, GitHub and repositories of AI-assisted mathematics) and found no new related work; the dataset~\cite{yang} was unchanged (still version 1.0.0, for orders $256$ and $384$), and that day the citation index OpenAlex refused to list the works citing~\cite{ollis16} (rate limit). Claude Code searched again on 22 November 2026 (05:39--05:51~UTC), when the OpenAlex lists of citing works refused on 21 November were retrieved, and found nothing new. This
reports the coverage of the searches, not a guarantee that they were complete.
The author is not a professional mathematician. The author has read this paper and, using a side-by-side table, compared its Lean statement with the main theorem, Theorem~\ref{thm:main}. No human
mathematician has reviewed this paper. The correctness of
Theorem~\ref{thm:main} rests on the Lean proof and on the agreement of the
Lean statement with the definition of a terrace in~\cite{ollis16,survey}; the
written proof is given for the reader, and the consequences in
Section~\ref{sec:context} rest on the cited theorems. AI systems are not
authors of this paper; the author takes full responsibility for its content.

\medskip
\noindent Thanks are due to M.~A. Ollis for the survey and the paper on
which the definitions and the context of this note rest, and to R.~A. Bailey
for the conjecture.

{\small
\begin{thebibliography}{9}
\setlength{\itemsep}{1pt}
\bibitem{ai93a} B.~A. Anderson and E.~C. Ihrig, Every finite solvable group
with a unique element of order two, except the quaternion group, has a
symmetric sequencing, \emph{J. Combin. Des.} 1 (1993), no.~1, 3--14,
\url{https://doi.org/10.1002/jcd.3180010103}. Not consulted; cited after~\cite{survey}.
\bibitem{ai93b} B.~A. Anderson and E.~C. Ihrig, Symmetric sequencings of
non-solvable groups, \emph{Congr. Numer.} 93 (1993), 73--82. Not consulted;
cited after~\cite{ollis16,survey}.
\bibitem{bailey} R.~A. Bailey, Quasi-complete Latin squares: construction and
randomization, \emph{J. R. Stat. Soc. Ser. B} 46 (1984), 323--334. Not
consulted; cited after~\cite{ollis16,survey}.
\bibitem{repo} H.~Dong, \emph{ai4math-results}, repository,
\url{https://github.com/Horace-Maxwell/ai4math-results}, version~\repoversion{}
(2026); archived at Zenodo under the concept DOI
\url{https://doi.org/10.5281/zenodo.22970254}.
\bibitem{mathlib}
The mathlib Community, The Lean mathematical library, in \emph{Proc.\ 9th ACM
SIGPLAN Int.\ Conf.\ Certified Programs and Proofs} (CPP '20), ACM, 2020,
\url{https://doi.org/10.1145/3372885.3373824}.
\bibitem{mp} A.~Müyesser and A.~Pokrovskiy, A random Hall--Paige conjecture,
\emph{Invent. Math.} 240 (2025), no.~3, 779--867,
\url{https://doi.org/10.1007/s00222-025-01328-x}; arXiv:2204.09666v3 (2025).
Theorem numbers refer to the arXiv version.
\bibitem{ollis16} M.~A. Ollis, Terraces for small groups, \emph{J. Combin.
Math. Combin. Comput.} 108 (2019), 231--244; arXiv:1603.01496v1 (2016).
Section, theorem and page numbers refer to the arXiv version.
\bibitem{survey} M.~A. Ollis, Sequenceable groups and related topics,
\emph{Electron. J. Combin.} 32(1) (2025), Dynamic Survey DS10v3 (version~3,
14 February 2025), \url{https://doi.org/10.37236/30}.
\bibitem{yang} H.~Yang, An explicit terrace for every non-abelian group of
order 256 and 384, dataset, version 1.0.0, Zenodo, 16 September 2026,
\url{https://doi.org/10.5281/zenodo.22785373}; description at
\url{https://severinvisionary.github.io/bailey-terraces-256-384/}.
\end{thebibliography}}
\end{document}
